% =============================================================================
% Abstract in italiano e in inglese, su due pagine separate, come nel template di
% riferimento (che ha Abstract + Résumé). Una sola voce nell'indice per entrambi.
% Testi da guideline/Intro&Abstract.docx, riportati parola per parola.
% =============================================================================
\cleardoublepage
\chapter*{Abstract}
~\newline~\newline~
\addcontentsline{toc}{chapter}{Abstract (Italiano/English)}
\pagestyle{plain}
\thispagestyle{plain}

\setlength{\parskip}{1em}

Il presente lavoro esplora le potenzialità informative del Baby-FE nell'analisi dei processi di regolazione in bambini che presentano condizioni di rischio evolutivo. All'interno di un campione di 66 bambini sono stati approfonditi otto casi, selezionati a priori sulla base delle informazioni anamnestiche raccolte dai genitori. Le prestazioni al Baby-FE e le osservazioni qualitative raccolte durante la somministrazione sono state integrate con le osservazioni comportamentali della Bayley-III e con le valutazioni delle insegnanti mediante EEFQ. I risultati evidenziano una marcata eterogeneità tra i casi e mostrano che, in prossimità del pavimento, punteggi complessivi molto bassi possono corrispondere a profili qualitativamente differenti. In tali condizioni, l'esame delle singole prestazioni e la distinzione tra prove non valide e prestazioni non corrette forniscono informazioni che il punteggio complessivo restituisce solo parzialmente. Nel complesso, il Baby-FE fornisce indicazioni utili, pur evidenziando limiti di sensibilità nelle fasce più basse di prestazione, che richiedono cautela interpretativa e ulteriori approfondimenti.

\setlength{\parskip}{0em}

\begin{otherlanguage}{english}
\cleardoublepage
\chapter*{Abstract}
~\newline~\newline~
\thispagestyle{plain}

\setlength{\parskip}{1em}

This study explores the informative potential of the Baby-FE for the analysis of regulatory processes in children presenting developmental risk factors. Within a sample of 66 children, eight cases were examined in greater depth and selected a priori on the basis of anamnestic information provided by their parents. Baby-FE performance and qualitative observations collected during administration were integrated with behavioral observations from the Bayley-III and teacher ratings obtained through the EEFQ. The results reveal marked heterogeneity across cases and show that, near the floor of the scale, very low total scores may correspond to qualitatively different profiles. Under these conditions, the examination of performance on individual tasks and the distinction between invalid trials and incorrect responses provide information that is only partially captured by the total score. Overall, the Baby-FE provides useful information, while also showing limitations in sensitivity at the lower end of the performance range, which call for cautious interpretation and further investigation.

\setlength{\parskip}{0em}
\end{otherlanguage}
